% Propietario conservado: /Users/ruben/Documents/New project/output/REVISION_ARGUMENTAL_APERTURAS_CIERRES_20260915/VI/delivery/MANUSCRIPT_VI_ES_EN_COMPLETE/source_es/sections/02_particula_persistencia.tex
\section{The particle as a class of persistent sections}
\label{vi:vi:sec:particula}

The material construction begins with the state produced by APP, TRIT, and
TPK. APP supplies the positions and the two arithmetic operations;
TRIT distinguishes the regime and orientation of their transitions; TPK
transports this information together with the quotients, boundary incidences,
and memory. As the state is prolonged, a finite evaluation
remains linked to the evaluations preceding it. The notion of a
particle rests on this compatibility and on its material representation.
Its mass will be an observable of the section thus constructed.

This sequence determines the meaning of the three objects to be
used: the genealogical state preserves the complete extension; the route
records its observable history; the material representation makes
that history act on a state space and permits the evaluation of observables.
The distinction is essential even when the three objects share
the same electronic label or the same atlas coordinate.

\subsection{Extension, observable history, and compatible section}
\label{vi:vi:subsec:extension-persistente}

Let $\mathcal S_m$ be the admissible states surviving through depth
$m$, obtained by TPK prolongation, and let
$\rho_{n,m}:\mathcal S_n\to\mathcal S_m$, for $n\geq m$, be their
restriction maps. Compatibility of refinement means
\begin{equation}
 \rho_{m,m}=I,\qquad
 \rho_{n,m}\rho_{p,n}=\rho_{p,m}\quad(p\geq n\geq m).
 \label{vi:vi:eq:sistema-estados}
\end{equation}
A surviving extension is an element
\begin{equation}
 x=(x_m)_{m\geq0}\in\Omega^{\mathrm{surv}}_{\mathrm{HMT}}
 :=\varprojlim_m\mathcal S_m,
 \qquad \rho_{n,m}(x_n)=x_m.
 \label{vi:vi:eq:extension-superviviente}
\end{equation}
The existence and admissibility of these states arise from the
construction of the formal kernel. The inverse-limit expression describes
the compatibilities preserved by a state already generated.

Let $\mathcal R_m$ be the space of observable histories at depth $m$.
The route reader $\operatorname{sh}_m:\mathcal S_m\to\mathcal R_m$
and the restriction $\operatorname{tr}_{n,m}$ satisfy
\begin{equation}
 \operatorname{tr}_{n,m}\operatorname{sh}_n
 =\operatorname{sh}_m\rho_{n,m}.
 \label{vi:vi:eq:naturalidad-ruta}
\end{equation}
Thus, the reader transmits the restrictions of the state to its histories.
The route of $x$ is the family
$\gamma_x=(\gamma_x^{[m]})_m$, where
$\gamma_x^{[m]}=\operatorname{sh}_m(x_m)$.

\begin{proposition}[Compatibility of the published history]
\label{vi:vi:prop:historia-compatible}
For every extension \eqref{vi:vi:eq:extension-superviviente}, its finite
histories satisfy
$\operatorname{tr}_{n,m}(\gamma_x^{[n]})=\gamma_x^{[m]}$.
\end{proposition}
\begin{proof}
Applying \eqref{vi:vi:eq:naturalidad-ruta} to the state $x_n$ gives
\[
 \operatorname{tr}_{n,m}(\gamma_x^{[n]})
 =\operatorname{sh}_m(\rho_{n,m}x_n)
 =\operatorname{sh}_m(x_m)=\gamma_x^{[m]}.
\]
The second equality uses precisely the inverse-limit condition.
\end{proof}

The proposition establishes descent from the extension to the history.
Its content does not include injectivity of the reader: two extensions
may publish the same finite history and retain different memories.
Therefore, when reusing a route, its source state is retained,
instead of reconstructing it from a digit or a visible word.

Over each history $\gamma_x^{[m]}$, consider the discrete bundle
$\mathcal B_{\gamma_x^{[m]}}$ of material states of the declared
representation. Let
$\mathscr E_m(x)=\Gamma(\mathcal B_{\gamma_x^{[m]}})$ be its spaces
of sections, with restrictions
$p_{n,m}:\mathscr E_n(x)\to\mathscr E_m(x)$ compatible with those
of the state. A persistent section is a family
\begin{equation}
 s=(s_m)_{m\geq m_0},\qquad s_m\in\mathscr E_m(x),\qquad
 p_{n,m}s_n=s_m\quad(n\geq m\geq m_0).
 \label{vi:vi:eq:seccion-persistente}
\end{equation}
The initial depth $m_0$ allows a state to be described from any
cut beyond which its representation is defined. Prior information
entering its observables remains in the record
transported by the extension.

\subsection{Identity under refinement and equivalence of representations}
\label{vi:vi:subsec:identidad-material}

Fix the internal symmetry groups $G_m$ acting on the material
fibres and preserving the declared observables. Two persistent
sections $s$ and $s'$ represent the same material identity when,
from some level $m_1$ onward, there exists a family $g_m\in G_m$ such that
\begin{equation}
 s'_m=g_ms_m,\qquad
 p_{n,m}g_n=g_mp_{n,m},
 \label{vi:vi:eq:equivalencia-persistente}
\end{equation}
and the observable characters of the genealogical record coincide. The
admissible symmetries are fixed before forming the quotient; their action
transports sheet, orientation, memory, and representation according to
the corresponding laws.

\begin{proposition}[Compatible equivalence of sections]
\label{vi:vi:prop:equivalencia-persistente}
The relation \eqref{vi:vi:eq:equivalencia-persistente}, on families carrying
the same types of representation and observables, is an equivalence
relation. Its classes are independent of removing a finite
number of initial levels, provided their transported genealogical
information is preserved.
\end{proposition}
\begin{proof}
The identity of each $G_m$ gives reflexivity. The inverse of a compatible
family is again compatible: from
$p_{n,m}g_n=g_mp_{n,m}$ one deduces
$p_{n,m}g_n^{-1}=g_m^{-1}p_{n,m}$. This gives symmetry.
If $s'=gs$ and $s''=hs'$, from the greater of the two initial levels
we have $s''_m=h_mg_ms_m$ and
$p_{n,m}h_ng_n=h_mg_mp_{n,m}$. Composition also preserves the
observables and proves transitivity. All these relations are formulated
on a common tail; changing the starting point within that tail
preserves the class. The final assertion concerns a change of
presentation, not erasure of state memory.
\end{proof}

\begin{definition}[Persistent material particle]
\label{vi:vi:def:particula}
An HMT--MD particle is a class $[s]_{\mathrm{pers}}$ of sections
\eqref{vi:vi:eq:seccion-persistente} under
\eqref{vi:vi:eq:equivalencia-persistente}, accompanied by its extension and
provenance record, admitting an irreducible material
representation or a minimal invariant block and a stable spectral output
under the specified coupling.
\end{definition}

The definition gathers conditions at different levels. Prolongation
and the record belong to the HMT state. Representation and
coupling belong to its material realisation. Spectral
stability is predicated of the operator and domain used by that
coupling. This organisation permits identity,
charge, spin, and mass to be studied together without turning one particular observable into the definition
of all the others.

The nonadic return illustrates the distinction between identity and instantaneous
state. For the phase reader $g$ and memory record
$\mathcal B$, the following can coexist:
\begin{equation}
 g(k+9)=g(k),\qquad
 p_{k+9,k}(s_{k+9})=s_k,\qquad
 \mathcal B_{k+9}\ne\mathcal B_k.
 \label{vi:vi:eq:persistencia-fase-memoria}
\end{equation}
The first equality is periodic; the second expresses persistence; the
third records advancement. The phase return belongs to the history
of the same material identity, while the state preserves the fact that
an additional circuit has been completed.

\subsection{The species as a fibre: scope of the finite atlas}
\label{vi:vi:subsec:especie-fibra}

The atlas constructed by the APP--TRIT--TPK operations has support
$\mathcal C_{81}=\{1,\ldots,9\}^2$. Its four orientations per
cell form the domain of $324$ oriented routes. The register
map $\mathcal R$ distinguishes $56$ route families, while the
internal classifier
\begin{equation}
 q:\mathcal C_{81}\longrightarrow\Sigma_{13}
 \label{vi:vi:eq:clasificador-especies}
\end{equation}
brings together $13$ family and level types. These censuses and their rules are
developed in the body devoted to the atlas; here we specify which object
they transmit to the definition of a particle.

For $y\in\Sigma_{13}$, the fibre $q^{-1}(y)$ preserves all its
cells and records. Once the internal ordering of each fibre of
$\mathcal R$ is fixed, let $k(c)$ be the rank of the cell within it. Its
finite presentation is
\begin{equation}
 \mathfrak S_y^{\mathrm{atlas}}
 =\{(\mathcal R(c),k(c)):q(c)=y\}.
 \label{vi:vi:eq:multiseccion-atlas}
\end{equation}
The rank is added to distinguish cells having the same family
register; it does not select a mass. The pair $(\mathcal R(c),k(c))$
recovers the cell within the ordered chart because the second datum
identifies its position in the fibre of the first.

If $\pi_{81}:\Omega^{\mathrm{surv}}_{\mathrm{HMT}}\to
\mathcal C_{81}$ is the cellular reader, the lift of the species is
\begin{equation}
 \widetilde{\mathfrak S}_y
 =\{x\in\Omega^{\mathrm{surv}}_{\mathrm{HMT}}:
                  q(\pi_{81}x)=y\}.
 \label{vi:vi:eq:multiseccion-enriquecida}
\end{equation}
This preimage contains extensions, whereas
\eqref{vi:vi:eq:multiseccion-atlas} contains finite presentations. The
construction and domain of the reader $\pi_{81}$ form part of the passage
between them. A particular particle is realised as a persistent section
within the corresponding enriched multisection.

\begin{theorem}[Central fixed point and equivariant selection]
\label{vi:vi:thm:selector-central}
Let $\rho_c(r,s)=(10-r,10-s)$ be cellular inversion. In the
$\rho_c$-stable fibres of the classifier \eqref{vi:vi:eq:clasificador-especies},
with trivial action on the label $y$, an equivariant point
selection can take values only in cell $(5,5)$. The electronic fibre
at level zero contains this unique fixed point; no other
fibre admits such a selection.
\end{theorem}
\begin{proof}
Solving $\rho_c(r,s)=(r,s)$ gives $r=s=5$. If $a(y)\in q^{-1}(y)$
is an equivariant selection and $y$ is fixed, then
$a(y)=a(\rho_c y)=\rho_c a(y)$. Its value must be the fixed point
$(5,5)$. The atlas places that cell in the electronic fibre at level
zero, whose cardinality is one. For any other fibre, the selection
condition would produce a fixed point not belonging to it.
\end{proof}

The consequence is constructive: the other species are presented through
orbits or multisections, and a realization requiring a particular
line must declare its representation and coupling. The
electron provides the central rank-one case without turning all
the families into copies of that case. Moreover, the cell $(5,5)$, the raw
central signature, the regional signature $555555$ and the relative origin of an
action line retain their own domains. Uniqueness of the point
does not identify these registers with one another.

\subsection{Construction of the mass observable from the material register}
\label{vi:vi:subsec:particula-observable}

The chain used by the mass observable preserves the order
\begin{equation}
 \text{extension}\longrightarrow\text{route}\longrightarrow
 \text{record}\longrightarrow\text{integer signature}\longrightarrow
 \text{character}\longrightarrow\text{towers}\longrightarrow
 \text{material operator}.
 \label{vi:vi:eq:cadena-material}
\end{equation}
Each arrow retains the data needed for the next. The signature is
a reading of the register; the character evaluates that signature with the
internal coefficients already constructed; the towers incorporate its
refinement; the operator finally acts in the representation of the
fibre selected by the coupling. Its output may be a scalar,
a multiplet or a spectral distribution, depending on that representation.
Metrological comparison takes place after this construction.

A species is described by charge, spin representation, colour,
flavour, generation, route family and material operator. These data
explain why two preparations with coincident mass may
retain different information, and why the same identity may
publish different observables under different couplings.
Identity resides in the persistent class and its admissible transports;
measurement is a reading of that identity in the declared codomain.

The same hierarchy separates a persistent extension from a section with a
finite horizon. If $\operatorname{Ext}_L(s)\ne\varnothing$ and
$\operatorname{Ext}_{L+1}(s)=\varnothing$, the horizon $L$ records
an obstruction to prolongation. An intermediate participation of this
kind in a composition does not by itself create a new asymptotic
species. Similarly, a resonance requires the survival
operator and its pole, as well as the identity of the section
it transports. The distinction between persistence,
finite prolongation, and spectral evaluation is thus preserved.
